\section{A Two-Level Framework for Noise Trajectory Search}
\label{sec:framework}

We cast noise trajectory search as a two-level problem: a general transition abstraction covers all diffusion-family models, on top of which a local search operator and a global scheduler serve as the two components.

\paragraph{General transition abstraction.}
Every diffusion-family model generates a sample by traversing $T$ discrete timesteps $t_T > t_{T-1} > \dots > t_1$. At each step the sampler applies
\begin{equation}
  x_{t-1} = F_\theta(x_t, t, c, \varepsilon_t),
  \label{eq:transition}
\end{equation}
where $x_t$ is the latent at timestep $t$, $c$ is the conditioning signal, and $\varepsilon_t$ is a noise variable injecting stochasticity. The final step yields $x_0$, which the verifier $v(x_0, c)$ scores. The transition is deliberately model-agnostic: for stochastic samplers (DDPM, SDE) $\varepsilon_t$ is the injected Gaussian noise; for deterministic samplers (flow-matching ODE solvers~\citep{lipman2023flow,liu2023flow}) stochasticity can be injected by adding controlled random perturbations~\citep{dockhorn2024stochastic,kim2025rbf,black2023training,fan2024reinforcement}.

\begin{figure}[t]
\centering
\includegraphics[width=0.92\linewidth]{figures/fig_architecture.pdf}
\caption{Two-level view of noise trajectory search. The \emph{global scheduler}~$\mathcal{G}$ (top) receives budget~$B$ and distributes per-step budgets $\{K_t\}$ across the trajectory. At each timestep, the \emph{local search operator}~$\mathcal{L}_t$ (bottom) evaluates $K_t$ noise candidates and selects the best. Bars depict the non-uniform allocation: timesteps where replacing the noise candidate yields larger quality gains receive bigger budgets.}
\label{fig:architecture}
\end{figure}

\paragraph{Local search operator.}
At timestep $t$, given $(x_t, c, v, F_\theta)$ and a per-step budget, the local operator
$\mathcal{L}_t : (x_t, c, v, F_\theta, K_t) \longrightarrow (x_{t-1}, \varepsilon_t)$
samples $K$ noise candidates $\{\varepsilon_t^{(k)}\}_{k=1}^K$, applies the base sampler to each to obtain a clean-data prediction
\begin{equation}
  x_{t-1}^{(k)} = F_\theta(x_t, t, c, \varepsilon_t^{(k)}),
  \qquad
  \hat{x}_0^{(k)} = D_\theta(x_{t-1}^{(k)}, t-1, c),
  \label{eq:x0_prediction}
\end{equation}
where $D_\theta$ is a one-step denoising formula~\citep{song2021scorebased}, and selects the candidate that maximizes the verifier score $s^{(k)} = v(\hat{x}_0^{(k)}, c)$. Each candidate costs one NFE. Many existing search procedures (zero-order perturbation, $\epsilon$-greedy exploration, particle-based refinement) instantiate this operator; we treat $\mathcal{L}_t$ abstractly as a black-box mapping from $x_t$ to $x_{t-1}$ given a per-step budget.

\paragraph{Global scheduler.}
The local operator decides \emph{how} to search at a given step but not \emph{where along the trajectory} to invest computation. A global scheduler
$\mathcal{G} : \bigl(x_T, c, v, F_\theta, B, \{\mathcal{L}_t\}_{t=1}^T\bigr) \longrightarrow \bigl(x_0, \{x_t\}_{t=1}^T\bigr)$
walks timesteps from $T$ to $1$ and at each step either invokes $\mathcal{L}_t$ (spending part of the remaining budget) or applies a single base step. This decision must be made \emph{online}: at step $t$ we cannot predict how effective search will be at later timesteps, yet computation spent now is lost forever.

\paragraph{Existing methods as special cases.}
This view unifies many prior strategies: plain diffusion sampling ($\mathcal{L}_t$ trivial); best-of-$N$ random search ($\mathcal{L}_t$ at the initial noise only); noise-trajectory and tree-style search (specific $\mathcal{L}_t$ paired with hard-coded scheduling). \cref{sec:mdp} presents a more general MDP formulation that subsumes these and accommodates richer action spaces (multiple verifier types, backtracking).
